\section{Related Work}
\label{sec:related}
\textbf{Multimodal Large Language Models}  
Recent advancements in Vision-Language Models (VLMs) have demonstrated exceptional multimodal understanding. Proprietary models~\cite{hurst2024gpt4o, anthropic2024claude3.5, wuevaluating, Google2023gemini} and open-source alternatives~\cite{Qwen2.5-VL, chen2024internvl, li2024llavaov, abdin2024phi, beyer2024paligemma} have set new benchmarks in visual question answering (VQA), image captioning, and multimodal dialogue through large-scale pretraining on image-text pairs. 
Reasoning-enhanced models like GPT-o1~\cite{OpenAI2024o1}, DeepSeek-R1~\cite{guo2025deepseek}, and Kimi-1.5~\cite{team2025kimi} show that post-training reinforcement learning (RL) can significantly improve mathematical and coding abilities. RL-based reasoning MLLMs~\cite{tan2025reason, huang2025vision} also excel in multimodal reasoning tasks. 
However, transferring these capabilities to embodied intelligence systems is a challenge. Works like EmbodiedGPT~\cite{mu2023embodiedgpt} and RoboBrain~\cite{ji2025robobrain} explore integrating visual-language understanding with robot-specific skills, including long-horizon task planning and trajectory synthesis. Extending reasoning-enhanced MLLMs to embodied scenarios~\cite{tan2025reason,  zhang2025embodied} represents a promising research frontier.


\textbf{Vision-Language-Action (VLA) Models}  
Building on the capabilities of VLMs~\cite{ chen2024internvl, li2024llavaov, beyer2024paligemma, li2024foundation}, researchers have developed vision-language-action (VLA) models for robotic manipulation tasks. Current VLAs are categorized as follows:
End-to-End VLAs directly map visual-textual inputs to actions using unified architectures, employing regression-based~\cite{li2023vision, brohan2023rt, kim2024openvla, hao2025tla, liu2024robomamba}, diffusion-based~\cite{liu2024rdt, black2024pi_0, team2024octo}, and hybrid approaches~\cite{liu2025hybridvla}. While effective for short-term tasks, they struggle with long-horizon planning and generalization.
Hierarchical VLAs address these limitations by using model-level~\cite{liu2024self, Helix2024, team2025gemini_robo, bjorck2025gr00t} or task-level hierarchies~\cite{shi2025hi, pi05} to decompose long-horizon tasks into subtasks. However, they still face challenges in cross-embodiment compatibility and multi-agent coordination.
To tackle these issues, we introduce RoboOS, the first open-source hierarchical embodied framework that facilitates cross-embodiment and multi-agent collaboration, achieving structural decoupling while maintaining functional synergy between MLLMs and VLAs.

\textbf{Multi-Robot Collaboration}  
Multi-robot collaboration (MRC) has been extensively studied for various applications~\cite{amazon_robotics_2025, mandi2024roco, an2023multi, liu2024coherent,HAO2025103018}. Research focuses on coordination, communication, and task allocation to enhance efficiency~\cite{rizk2019cooperative, fierro2018multi}. MRC shows promise in automated warehousing~\cite{agrawal2023rtaw}, search and rescue~\cite{guo2023cross}, and environmental monitoring~\cite{edwards2023collaborative}. Learning methods like reinforcement and imitation learning further improve MRC~\cite{patino2023learning, liu2023guide, pereida2018data}.
Despite advancements, MRC faces challenges in cross-embodiment adaptation, task allocation, and dynamic planning, impacting real-time coordination. This paper introduces RoboOS, a hierarchical architecture that addresses these issues and enhances edge-cloud communication for scalable deployment.


